% !TeX root = ../main.tex
\chapter{选题依据}
\label{chap:background}

\section{研究背景}
本节填写研究问题所处的应用背景、已有技术条件及现实需求。建议从具体场景出发，
说明现有方法遇到的限制，再引出需要研究的问题。涉及现状、性能或趋势的判断应有可靠文献支持。

\section{研究意义}
分别说明本课题可能提供的理论认识、方法改进或工程价值，明确这些价值对应的对象和使用条件。
开题阶段应描述拟研究的内容与预期贡献，避免把尚未开展的工作写成已经取得的成果。

\section{研究问题与范围}
用一段话给出问题定义，说明研究对象、输入输出、约束条件以及不在本课题范围内的内容。
随后列出两到三个可检验的问题，例如：在给定约束下，如何改善目标指标，以及如何验证改进来自所提出的方法。
